\label{sec:ch17-business-migration}
\subsection*{冻结旧权威的闭合前缀}

现在将 b1 从 G1 迁到 G2。迁移前 G1 持有归属代次 1，G2 尚无激活的 b1。迁移服务取得唯一迁移标识 M1，向 G1 提交 freeze 命令。freeze 在 G1 的业务日志中排序，应用到某个位置后把 b1 标为 \code{frozen}，并记录 M1。此前已经提交的变化属于迁移图像，此后的业务修改按冻结规则退出。

冻结覆盖所有改变图像的路径：新受理、原结果观察、会话注册与关闭、工作者 claim、退款受理与结果、发布标记、交接回执与履约进度，以及该分区拥有的库存和事务状态。只停止新 $S$ 绑定，旧工作者还可能回填平台结果，或把退款由 \code{pending} 改为 \code{succeeded}，图像就会在复制期间变化。完整冻结让导出内容对应稳定业务前缀。

\CSeventeenFigure{ch17-13-handoff}{源冻结、图像导入、目录条件发布和目标激活依次交接归属}{fig:ch17-handoff}

冻结时若存在尚未闭合的 prepared 库存责任，系统采用明确的迁移准入规则。本版在有 prepared 参与者或未完成协调者事务时拒绝冻结，等待对应决定完成后再发起迁移。另一种系统可以连同锁、参与者身份、协调者决定位置和恢复路由一起搬迁，但那要求各参与者识别新的责任位置。此处首版实现用实际拒绝固定边界，待提交责任保持在原组继续解决。

\CSeventeenFigure{ch17-24-freeze-gate}{prepared 责任未闭合时保留原权威，满足准入后才冻结}{fig:ch17-freeze-gate}

冻结后其他逻辑分区继续工作。G1 可以还承载其他分区，G2 和 G3 也继续服务自己已有分区。b1 的暂停限于其修改路径及确实依赖它的跨分区协作。入口可以提示稍后按原 $S$ 恢复，后台也可以保存待处理工作；恢复后的受理判断继续引用原输入。

\subsection*{图像导入与目录发布}

G1 导出冻结图像及摘要，M1、源归属代次和业务状态一起进入迁移记录。图像包含全部已受理操作和完整未完成责任。迁移服务持久保存这份记录，自身在任何阶段退出后，新实例可以重新读取源组冻结状态，确认重复 freeze 得到同一图像。

G2 校验图像格式、摘要和内部关联，再以更高归属代次 2 导入 b1。导入后先处于 \code{inactive}。此时 G2 已有数据，业务入口、支付 claim 和结果观察暂时不能开始新的修改。数据准备与服务资格分别持久保存，目录交接便有确定的中间状态。

迁移服务在控制面执行条件更新：当前目录等于 G1/1 时，才改为 G2/2，并写入 M1 和图像摘要。另一迁移者若先改变归属，本次比较就失败。服务读取已提交目录，核对 G2、代次 2、M1 和摘要全部相符，才沿受控管理路径激活 G2 的 b1。

\begin{BookCode}
{
  "bucket": "b1",
  "owner": "G2",
  "epoch": 2,
  "migration": "M1",
  "source_owner": "G1",
  "source_epoch": 1,
  "image_format": 1,
  "image_sha256":
    "9239bb141c68864ab7dc6b5b7b810126c9b9ceefcfb79d474456a5d28c9e0baf"
}
\end{BookCode}

这份记录发布迁移关系，由 64 个十六进制字符组成的 SHA-256 摘要标记所导入的冻结图像。源组保留 frozen 副本，目标根据已提交目录取得激活依据。管理接口属于受控迁移服务的信任边界，任意调用者若能绕过该路径直接激活目标，就会破坏交接次序。管理资格和业务访问资格分别保护，激活入口核验的依据也保存到目标状态中。

激活响应也可能丢失。恢复者重读目录与目标模式，若 M1、摘要和代次匹配，就完成同一交接。目标的原子激活入口只把匹配的 inactive 改为 active；当前已经 active 且同一交接匹配时返回已完成事实。目标后来为 M2 冻结，或已保存更高归属代次时，迟到的 M1 激活进入冲突，保持当前模式。一次目录发布结果不确定时，重新创建空 b1 并换迁移标识，会使新旧图像失去共同依据。各阶段重复操作都比较原迁移标识和内容，返回同一阶段已经成立的事实。

\subsection*{迟到请求与原平台操作}

迁移过程中入口 A 一直缓存 G1/1。切换后它将 Alice 原 $S$ 发给 G1，源组按 frozen 模式拒绝修改。A 刷新目录，再把同一 $S$ 和原输入送到 G2/2。原操作已有结果时，G2 返回图像中的原结果；操作处于 pending 时，恢复者沿原 C、原范围和 \code{pay-C} 查询或调用平台。

旧工作者 B 也可能在冻结前已发出付款 HTTP。平台在冻结后完成 X1，B 随后向 G1 回填证据，源组冻结规则使修改退出。新工作者从 G2 待办中取得同一平台键，得到 X1 原证据，再用 G2 当前工作代次提交观察。迟到通信与新的恢复扫描汇合到同一平台操作。

\CSeventeenFigure{ch17-14-late-work}{旧路由在源组退出，新位置用原平台身份取得同一付款证据}{fig:ch17-late-work}

这里分别检查两个状态：源冻结图像在导出后保持不变，目标从同一图像出发继续前进。目标后来会增加付款结果或发布标记，核对时比较原 $S$、原输入、P/C/T、平台范围、键和已保存原结果。目标全部字节与冻结图像的差异，可能正是交接后合法发生的新历史。

若只复制订单表，G2 可能看不到尚未形成订单的 pending 付款。再次受理同 $S$ 时，它就可能分配另一个 C 并使用另一平台键。两个组各自的本地唯一约束都成立，平台却接收两项不同操作。闭合图像带走未决责任，归属交接让后续修改在一个有资格的位置继续。

最后把几种移动放到各自层次。etcd 成员更改改变协调集群的投票与承载，购买组的 Raft 成员更改改变该组副本。本节交接包含会话、身份、结果和待办的完整业务分区。下一节的 TiKV Peer 调度和 Region split/merge 则处理底层区间与复制布局，第十八章的 Redis Cluster 槽迁移处理槽中键的所在节点。

\begin{table}[htbp]\centering\small
\begin{tabularx}{\linewidth}{@{}p{.22\linewidth}XX@{}}
\toprule
移动的对象 & 保持连续的状态 & 完成依据\\
\midrule
协调或存储复制组的成员 & 同一组的日志、快照与配置历史 & 新副本追赶及对应成员更改协议完成\\
Region 或槽中的数据 & 区间或槽内的键及产品规定元数据 & 产品的范围、成员、路由及导入导出步骤完成\\
完整购买业务分区 & 会话、$S$ 绑定、原结果、C/T、退款及未完成责任 & 源冻结、图像导入、目录发布与目标激活闭合\\
\bottomrule
\end{tabularx}
\caption{不同层次移动的对象与完成依据}\label{tab:ch17-moves}
\end{table}
\FloatBarrier

业务分区交接完成后，入口重新定位，恢复者继续处理原待办，底层复制组保护目标上的新历史。如果用通用分布式数据库承载这些业务表，数据库还会独立调整底层键空间。接下来看 TiKV 怎样移动 Peer、分裂和合并 Region，比较这一层的范围变化与刚完成的完整业务分区交接。
\FloatBarrier
